% !TeX root = ../../Book3.tex
% 纲要标题：## 第13章　滤过、Artin–Rees 与完备化
% 纲要位置：工作记录/计划纲要.md，第 2335 行。

\chapter{滤过、Artin–Rees 与完备化}
\label{b3:ch:13}
\BookChapterNumber{b3:ch:13}

\begin{chapterabstract}
本章把环与模的完备化作为理想幂商的逆极限来构造。Rees 环的有限性导出 Artin–Rees 引理和 Jacobson 根条件下的 Krull 交定理；随后研究有限生成模的完备化何时保持正合，并讨论完备环的 Noether 性、忠实平坦性与局部维数。
\end{chapterabstract}

\begin{learninggoals}
\begin{itemize}
\item 比较子模自身的理想进滤过与环境诱导的滤过，能指出 Artin–Rees 中有限性的作用。
\item 用相容余类构造完备化，解释逆极限一般不保持满射的原因。
\item 由完备化正合性和张量表示证明平坦性，并用剩余域检测忠实性。
\item 计算典型的幂级数完备化，并区分有限阶相等、形式相等和解析收敛。
\end{itemize}
\end{learninggoals}

整数的序列 \(1,1+2,1+2+4,\ldots\) 在通常绝对值下越走越远；模 \(2^n\) 看，第 \(n\) 步以后它却不再改变，且所得相容余数都等于 \(-1\) 的余数。这种收敛取决于理想的幂，而不是实数距离。把所有模 \(I^n\) 的相容答案合在一起，便是 \(I\) 进完备化的基本构造。

\ProseParagraphBreak

对一个模作同样的操作时，逐层商的正合性未必自动传到无限相容族；取子模后所诱导的滤过，也未必显然等于原子模自身的理想幂滤过。Artin–Rees 引理正是控制这两个滤过的差距。弄清这个有限性步骤，才能说明有限生成模的完备化何时保持正合，以及完备局部环为何仍服从 Noether 性。

\ProseParagraphBreak

本章使用\chapref{b3:ch:03}的 Noether 模与 Nakayama 引理、\secref{b3:sec:1104}的 Hilbert–Samuel 理论，以及\secref{b3:sec:1201}、\secref{b3:sec:1205}的平坦和忠实平坦判据。逆极限在本章从相容族定义。除非另有说明，环均交换且有单位。

\ProseParagraphBreak
Artin–Rees、有限生成模的张量表示及完备化性质可对照松村英之的《Commutative Ring Theory》\cite[\S8, Theorems 8.5--8.14, pp. 59--62]{Matsumura1989CommutativeRingTheory}。该书的证明将滤过与有限生成相联系，本章同时展开相容提升和逐阶消去的步骤，以说明完备性在何处真正起作用。


% !TeX root = ../../Book3.tex
% 纲要标题：### 13.1 理想进滤过
% 纲要位置：工作记录/计划纲要.md，第 2359 行。
% 纲要细目（写作范围）：
% * I 进拓扑与分离性
% * 模滤过和伴随分次
% * 相容余类族与滤过完备化；稳定滤过和理想幂滤过之间的自然同构
% * Rees 环与有限生成；理想幂滤过的 Rees 环及模模掉 \(I\) 后得到伴随分次对象，一般稳定滤过只保证在充分高次满足相应等式

\section{理想进滤过}
\label{b3:sec:1301}


多项式只含有限个系数，形式幂级数却要同时保存任意高阶的系数。把一个对象模掉 \(I,I^2,I^3,\ldots\)，就是逐次保留更多信息的代数方法。本节先说明“两个元素相差很小”的含义；这里的“小”由理想的幂衡量，不依赖绝对值。

\subsection{\texorpdfstring{\(I\)}{I} 进拓扑与分离性}
\label{b3:subsec:1301:01}

\begin{definition}\label{b3:def:adic-topology}
设 \(R\) 为交换环、\(I\subseteq R\) 为理想、\(M\) 为 \(R\) 模。以
\[
x+I^nM\qquad(x\in M,\ n\ge0)
\]
为邻域基，得到 \(M\) 的 \term[I-jin-tuopu]{\(I\) 进拓扑}[I-adic topology]：一个集合是开集，是指其中每个点都包含在某个仍位于该集合内的上述邻域中。若
\(\bigcap_{n\ge0}I^nM=0\)，则称这个拓扑\term[fenli-tuopu]{分离}[separated]。对 \(M\) 中的序列 \((x_j)\) 及元素 \(x\in M\)，若对每个 \(n\ge0\)，充分大的 \(j\) 都有 \(x_j-x\in I^nM\)，则称 \(x_j\) 收敛于 \(x\)；若对每个 \(n\ge0\)，充分大的 \(i,j\) 都有 \(x_i-x_j\in I^nM\)，则称 \((x_j)\) 为 Cauchy 序列\footnote{柯西（Augustin-Louis Cauchy，1789-1857），法国数学家，推动极限、收敛与复分析的严格化。}。
\end{definition}

由于 \(I^{n+1}M\subseteq I^nM\)，这些集合确实给出拓扑。商 \(M/I^nM\) 把相差 \(I^nM\) 中元素的两个值看作相同；自然映射 \(M/I^{n+1}M\to M/I^nM\) 则将较高一阶的精度降回第 \(n\) 阶。因此 \(x_j\to x\) 表示：对每个固定的 \(n\)，充分大的 \(j\) 都使 \(x_j\) 与 \(x\) 在 \(M/I^nM\) 中完全相同。加法与取负连续；\(R\) 自身的 \(I\) 进拓扑还使乘法连续，因为 \(I^n\) 是理想。

若拓扑分离，\(x\ne y\) 时就有某个 \(n\) 使 \(x-y\notin I^nM\)。此时两个邻域 \(x+I^nM\)、\(y+I^nM\) 不交，否则交点的两个表示会给出 \(x-y\in I^nM\)。反过来，若 \(0\ne x-y\in\bigcap_n I^nM\)，这两个点在每个商中都相同，不能用开邻域分开。因此分离性等价于 Hausdorff 性，也恰好表示各阶商合起来能够区分原来的不同元素。

\ProseParagraphBreak
设 \(k\) 为域。在 \(k[x]\) 中，\((x)\) 进收敛表示每个固定次数的系数最终不再改变。非零多项式不能被任意高次的 \(x\) 整除，所以此拓扑分离，收敛序列的极限至多有一个。但 \(1+x+\cdots+x^j\) 是 Cauchy 序列，任何极限的每个系数却都必须是一，不可能是多项式。它在 \(k[x]\) 内没有极限，在 \(k[[x]]\) 中才得到极限 \(\sum_{i\ge0}x^i\)。分离性保证已有极限的唯一性，补足这些极限则是另一个问题。另取 \(R=k\times k\)、\(I=k\times0\)，则 \(I^n=I\) 对所有 \(n\ge1\) 成立，拓扑不分离；第二个坐标以外的信息已经在第一个商中被永久遗失。

\begin{proposition}\label{b3:prop:adic-closure}
设 \(R\) 为交换环、\(I\subseteq R\) 为理想、\(M\) 为 \(R\) 模，并在 \(M\) 上取 \(I\) 进拓扑。子模 \(N\subseteq M\) 的闭包为
\[
\overline N=\bigcap_{n\ge0}(N+I^nM).
\]
特别地，\(N\) 闭当且仅当 \(M/N\) 的 \(I\) 进拓扑分离。对任意 \(R\) 模 \(P\)，每个 \(R\) 线性映射 \(f:M\to P\) 在两模的 \(I\) 进拓扑下都连续。
\end{proposition}
\begin{proof}
点 \(x\) 属于闭包，恰指对每个 \(n\)，邻域 \(x+I^nM\) 都与 \(N\) 相交。也就是可以选 \(y_n\in N\)，使 \(x-y_n\in I^nM\)：在任意给定精度下，都能用 \(N\) 中的元素逼近 \(x\)。这等价于 \(x\in N+I^nM\) 对每个 \(n\) 成立，得到闭包公式。又
\(I^n(M/N)=(I^nM+N)/N\)，得到闭性的刻画。若 \(f:M\rightarrow P\) 线性，则 \(f(I^nM)\subseteq I^nP\)，所以原点处连续；平移后处处连续。
\end{proof}

\subsection{模滤过与伴随分次}
\label{b3:subsec:1301:02}

\begin{definition}\label{b3:def:good-adic-filtration}
设 \(R\) 为交换环、\(M\) 为 \(R\) 模。下降的子模列
\(M=F^0M\supseteq F^1M\supseteq\cdots\) 称为一个\term[mo-lvguo]{模滤过}[module filtration]。给定理想 \(I\subseteq R\)，若 \(IF^nM\subseteq F^{n+1}M\) 对每个 \(n\ge0\) 成立，则称它为 \(I\) 滤过；若还有 \(IF^nM=F^{n+1}M\) 对充分大的 \(n\) 成立，则称它为\term[wending-lvguo]{稳定滤过}[stable filtration]。对 \(I\) 滤过，相应伴随分次模为
\[
\operatorname{gr}_F(M)=\bigoplus_{n\ge0}F^nM/F^{n+1}M.
\]
它是 \(\operatorname{gr}_I(R)\) 模，作用由代表元相乘定义。
\end{definition}

良定义需要同时检查两个代表元：若 \(a\in I^r\)、\(x\in F^nM\)，更换 \(a\) 增加的项属于 \(I^{r+1}F^nM\subseteq F^{n+r+1}M\)，更换 \(x\) 增加的项也属于此子模。\cref{b3:def:associated-graded-adic}中的滤过 \(F^nM=I^nM\) 是从第一步就稳定的特例。

\begin{definition}\label{b3:def:filtered-completion}
设 \(R\) 为交换环、\(M\) 为 \(R\) 模，\(M=F^0M\supseteq F^1M\supseteq\cdots\) 为下降滤过。对 \(n\ge1\)，记 \(q^F_n:M/F^{n+1}M\rightarrow M/F^nM\) 为自然商映射。该滤过的相容余类族组成 \(R\) 模
\[
\widehat M_F\coloneqq
\left\{\,(x_n)_{n\ge1}\in\prod_{n\ge1}M/F^nM
\mid q^F_n(x_{n+1})=x_n\text{ 对所有 }n\ge1\,\right\},
\]
加法和标量作用逐坐标进行。这个相容余类模也记为
\(\varprojlim_{n\ge1}M/F^nM\)。对理想 \(I\subseteq R\) 的幂滤过 \(F^nM=I^nM\)，将它记为 \(\widehat M_I\)。自然映射 \(M\rightarrow\widehat M_F\) 把 \(x\) 送到 \((x+F^nM)_n\)。
\end{definition}

伴随分次的第 \(n\) 片 \(F^nM/F^{n+1}M\) 记录这一层相对于下一层留下的信息；相容余类模则保留每个完整的截断 \(M/F^nM\)，以及高阶截断如何映到低阶截断。比较两套滤过时，若它们的精度只相差固定的有限阶数，要求达到任意高阶精度就会给出同一拓扑，并使两个相容余类模自然同构。

\begin{proposition}\label{b3:prop:stable-filtration-equivalent}
设 \(R\) 为交换环、\(I\subseteq R\) 为理想、\(M\) 为 \(R\) 模，\(F^nM\) 为从次数 \(c\ge0\) 起稳定的 \(I\) 滤过。则对 \(n\ge c\) 有
\[
F^nM=I^{n-c}F^cM,\qquad
I^nM\subseteq F^nM\subseteq I^{n-c}M.
\]
因而 \(F^nM\) 与 \(I^nM\) 定义同一拓扑，恒等映射 \(M\rightarrow M\) 还诱导自然同构
\[
\widehat M_I=\varprojlim_{n\ge1}M/I^nM
\xrightarrow{\ \sim\ }\varprojlim_{n\ge1}M/F^nM=\widehat M_F,
\]
它与从 \(M\) 到两个相容余类模的自然映射相容。
\end{proposition}
\begin{proof}
第一个等式由稳定性归纳得到。由 \(F^0M=M\) 及 \(IF^nM\subseteq F^{n+1}M\)，归纳得 \(I^nM\subseteq F^nM\) 对所有 \(n\ge0\) 成立；右侧包含由第一个等式得到。要使差值落入 \(F^nM\)，只需使它落入 \(I^nM\)；要使差值落入 \(I^nM\)，则使它落入 \(F^{n+c}M\) 即可。这就使两个邻域系相互共尾，拓扑相同。

对每个 \(n\ge1\)，包含 \(I^nM\subseteq F^nM\) 及 \(F^{n+c}M\subseteq I^nM\) 给出自然商映射
\[
\alpha_n:M/I^nM\longrightarrow M/F^nM,
\qquad
\beta_n:M/F^{n+c}M\longrightarrow M/I^nM.
\]
它们与各自的过渡商映射交换，所以逐坐标定义
\[
\Phi:\widehat M_I\longrightarrow\widehat M_F,
\quad (x_n)_n\longmapsto\bigl(\alpha_n(x_n)\bigr)_n,
\]
\[
\Psi:\widehat M_F\longrightarrow\widehat M_I,
\quad (y_n)_n\longmapsto\bigl(\beta_n(y_{n+c})\bigr)_n
\]
确实得到相容余类族。复合 \(\beta_n\alpha_{n+c}\) 是从 \(M/I^{n+c}M\) 到 \(M/I^nM\) 的过渡商映射，故 \((\Psi\Phi(x))_n=x_n\)。复合 \(\alpha_n\beta_n\) 是从 \(M/F^{n+c}M\) 到 \(M/F^nM\) 的过渡商映射，故 \((\Phi\Psi(y))_n=y_n\)。因此两映射互逆。

若将 \(c\) 增大，相容性保证 \(\Psi\) 的每个坐标不变。这些商映射都由 \(M\) 的恒等映射诱导，故所得到的同构与自然余类映射相容，也与保持滤过的模同态相容。
\end{proof}

伴随分次未必恢复原环。设 \(p\) 为素数，取 \(A=\mathbb Z/p^2\mathbb Z\)、\(B=\mathbb F_p[t]/(t^2)\)。它们的极大理想分别为 \((p)\) 与 \((t)\)，平方都为零，次数零的片都是 \(\mathbb F_p\)，次数一的片分别由 \(p\)、\(t\) 的类生成。因此两个伴随分次环都为 \(\mathbb F_p[T]/(T^2)\)。但在 \(A\) 中，\(p\) 个单位元之和为 \(p\ne0\)，落入下一层；在 \(B\) 中，\(p\) 个单位元之和已经为零。在伴随分次的次数零部分，这两个和都为零，因而看不到前一种跨到下一层的加法关系，原环的特征仍可不同。

\subsection{Rees 环与有限生成}
\label{b3:subsec:1301:03}

\begin{definition}\label{b3:def:rees-algebra-module}
设 \(R\) 为交换环、\(I\subseteq R\) 为理想、\(M\) 为 \(R\) 模，\(F^nM\) 为 \(I\) 滤过。引入形式变量 \(t\)，定义 \term[rees-huan]{Rees 环}[Rees algebra]与 Rees 模
\[
\mathcal R_I(R)=\bigoplus_{n\ge0}I^nt^n\subseteq R[t],\qquad
\mathcal R_F(M)=\bigoplus_{n\ge0}F^nM\,t^n.
\]
后者的作用为 \((at^r)(xt^n)=(ax)t^{r+n}\)。这只是给各阶加上标记；不要求原模中的元素本来具有次数。对幂滤过 \(F^nM=I^nM\)，将相应 Rees 模记为 \(\mathcal R_I(M)\)。
\symidx{\(\mathcal R_I(R)\)：理想的 Rees 环}
\symidx{\(\mathcal R_F(M)\)：滤过的 Rees 模}
\end{definition}

\cref{b3:lem:induced-filtration-bound}的证明已用这种构造同时处理全部次数。作为 Rees 环上的模，有限生成要求所有层都由有界次数内的有限组齐次元素产生；只知道每一层分别有限生成，还不能得到这样一组共同生成元。

\begin{proposition}\label{b3:prop:rees-finiteness}
设 \(R\) 为交换环、\(I\subseteq R\) 为理想、\(M\) 为 \(R\) 模。幂滤过的 Rees 环及模有自然的分次同构
\[
\mathcal R_I(R)/I\mathcal R_I(R)\cong\operatorname{gr}_I(R),
\qquad
\mathcal R_I(M)/I\mathcal R_I(M)\cong\operatorname{gr}_I(M),
\]
后一个同构的模作用与前一个环同构相容。若 \(R\) 为 Noether 环，则 \(\mathcal R_I(R)\) 为 Noether 环；若还有 \(M\) 有限生成，则对 \(M\) 上的 \(I\) 滤过，\(\mathcal R_F(M)\) 作为 \(\mathcal R_I(R)\) 模有限生成当且仅当滤过稳定。
\end{proposition}
\begin{proof}
理想 \(I\) 放在 Rees 环的次数零部分时，\(I\mathcal R_I(R)\) 的第 \(n\) 片为 \(I^{n+1}t^n\)，故商的第 \(n\) 片为 \(I^n/I^{n+1}\)。逐片把 \(at^n\) 的类送到 \(a+I^{n+1}\)，乘法与代表元相乘相容，就得到第一个同构。同样，\(I\mathcal R_I(M)\) 的第 \(n\) 片为 \(I^{n+1}M\,t^n\)，逐片取商得到第二个同构，且环对模的作用也由相同的代表元乘法给出。

以下设 \(R\) 为 Noether 环。写 \(I=(a_1,\ldots,a_s)\)，则 \(\mathcal R_I(R)=R[a_1t,\ldots,a_st]\)，由 Hilbert 基定理与商环的 Noether 性得到 Rees 环的 Noether 性。

再设 \(M\) 有限生成。若滤过从 \(c\) 起稳定，选 \(F^0M,\ldots,F^cM\) 各自的有限组生成元；这些子模有限生成是因为 \(M\) 为 Noether 模。把这些生成元放入相应次数后，由 \(F^nM=I^{n-c}F^cM\) 可知它们生成全部 Rees 模。

反向可取齐次生成元 \(x_jt^{d_j}\)：将任意有限组生成元分解成齐次部分，所得有限组仍生成。设 \(c=\max_jd_j\)，则
\[
F^nM=\sum_jI^{n-d_j}x_j\quad(n\ge c).
\]
于是 \(F^{n+1}M=IF^nM\) 对 \(n\ge c\) 成立。零模的情形直接成立。
\end{proof}

一般 \(I\) 滤过的 Rees 模模去 \(I\) 后，第 \(n\) 片是 \(F^nM/IF^nM\)，而伴随分次的第 \(n\) 片是 \(F^nM/F^{n+1}M\)。\(I\) 滤过只要求 \(IF^nM\subseteq F^{n+1}M\)，两种商未必相同。即使滤过稳定，也只能保证稳定以后的各次满足等式，不能据此将整个 Rees 商模认同为伴随分次模。

% !TeX root = ../../Book3.tex
% 纲要标题：### 13.2 Artin–Rees 引理
% 纲要位置：工作记录/计划纲要.md，第 2344 行。
% 纲要细目（写作范围）：
% * 子模滤过与诱导滤过
% * Artin–Rees 引理及证明
% * Krull 交定理与子模闭性：Jacobson 根条件及局部环特例

\section{Artin–Rees 引理}
\label{b3:sec:1302}


在子模 \(N\subseteq M\) 内，“\(I\) 的 \(n\) 次幂乘出的元素”组成 \(I^nN\)；“在 \(M\) 中已经达到第 \(n\) 阶的元素”组成 \(N\cap I^nM\)。这两者一般不相等。Artin–Rees 引理说明，在 Noether 条件下，它们的差异能由一个固定的次数偏移控制。

\subsection{子模滤过与诱导滤过}
\label{b3:subsec:1302:01}

设 \(k\) 为域，取 \(R=k[x]\)、\(I=(x)\)、\(M=R\)、\(N=x^rR\)，其中 \(r\ge1\) 为整数。直接计算得到
\[
I^nN=x^{n+r}R,\qquad N\cap I^nM=x^{\max(r,n)}R.
\]
后一滤过到第 \(r\) 阶仍等于整个 \(N\)，前一滤过已经缩小。因此子模的自身理想进滤过和由环境诱导的滤过需要分别处理。

\ProseParagraphBreak
令 \(F^nN=N\cap I^nM\)。它满足 \(F^0N=N\) 及 \(IF^nN\subseteq F^{n+1}N\)，且其 Rees 模是 \(\mathcal R_I(M)\) 的齐次子模。这个观察把滤过比较转化为一个有限生成问题。

\subsection{Artin–Rees 引理及其证明}
\label{b3:subsec:1302:02}

下述\term[artin-rees-yinli]{Artin–Rees 引理}[Artin--Rees lemma]给出诱导滤过最终稳定的精确等式。

\begin{theorem}[Artin–Rees]\label{b3:thm:artin-rees}
设 \(R\) 为交换 Noether 环、\(I\subseteq R\) 为理想、\(M\) 为有限生成 \(R\) 模、\(N\subseteq M\) 为子模。存在整数 \(c\ge0\)，使对所有 \(n\ge c\) 都有
\[
N\cap I^nM=I^{n-c}(N\cap I^cM).
\]
特别地，\(N\) 自身的 \(I\) 进拓扑等于 \(M\) 在 \(N\) 上诱导的拓扑。
\end{theorem}
\begin{proof}
因为 \(R\) 为 Noether 环、\(M\) 有限生成，子模 \(N\) 也是有限生成 \(R\) 模。由\cref{b3:prop:rees-finiteness}，\(\mathcal R_I(R)\) 为 Noether 环，\(\mathcal R_I(M)\) 为有限生成模，所以齐次子模
\[
\mathcal N=\bigoplus_{n\ge0}(N\cap I^nM)t^n
\]
有限生成。由于 \(N\) 有限生成，可再次应用该命题，得到滤过 \(N\cap I^nM\) 从某个次数 \(c\) 起稳定，逐次使用等式便得到所述公式。

具体看齐次生成元，可取 \(\mathcal N\) 的有限组生成元 \(z_jt^{d_j}\) 并增大 \(c\)，使 \(d_j\le c\)。对 \(n\ge c\)，齐次生成性给出
\[
N\cap I^nM=\sum_jI^{n-d_j}z_j.
\]
因为 \(z_j\in N\cap I^{d_j}M\)，有
\[
I^{n-d_j}z_j=I^{n-c}\bigl(I^{c-d_j}z_j\bigr),
\qquad I^{c-d_j}z_j\subseteq N\cap I^cM.
\]
故上面的和包含于 \(I^{n-c}(N\cap I^cM)\)。反向包含则由 \(N\) 为子模及 \(I^{n-c}I^cM\subseteq I^nM\) 得到，这证明所需的精确等式。

最后，
\[
I^nN\subseteq N\cap I^nM\subseteq I^{n-c}N\qquad(n\ge c).
\]
两个邻域系相互包含任意给定的足够小邻域，故拓扑相同。
\end{proof}

\cref{b3:lem:induced-filtration-bound}只需要上述两个包含关系。现在得到的等式还说明诱导滤过何时开始逐次由乘 \(I\) 生成。这里的 \(c\) 可以依赖 \(N\subseteq M\) 和 \(I\)，不能解释为对所有子模统一的常数；前面 \(N=x^rR\) 的例子中，最小的 \(c\) 就是 \(r\)。

\subsection{Krull 交定理与子模闭性}
\label{b3:subsec:1302:03}

\term[krull-jiao-dingli]{Krull 交定理}[Krull intersection theorem]以 Jacobson 根条件保证有限生成模的分离性。

\begin{theorem}[Krull 交定理]\label{b3:thm:krull-intersection}
设 \(R\) 为交换 Noether 环、\(M\) 为有限生成 \(R\) 模、\(I\subseteq R\) 为理想，\(J(R)\) 为 \(R\) 的 Jacobson 根。若 \(I\subseteq J(R)\)，则
\[
\bigcap_{n\ge0}I^nM=0.
\]
此时 \(M\) 的每个子模都闭。特别地，对 Noether 局部环 \((R,\mathfrak m)\)，每个有限生成模在极大理想进拓扑下分离。
\end{theorem}
\begin{proof}
令 \(N=\bigcap_nI^nM\)。它作为 Noether 模的子模有限生成。对 \(N\subseteq M\) 使用\cref{b3:thm:artin-rees}，在 \(n=c+1\) 时得到
\[
N=N\cap I^{c+1}M=I(N\cap I^cM)=IN.
\]
由\cref{b3:thm:nakayama}及 \(I\subseteq J(R)\) 得 \(N=0\)。对任意子模 \(L\)，商模 \(M/L\) 也有限生成，因而分离；\cref{b3:prop:adic-closure}说明 \(L\) 闭。
\end{proof}

有限性不能删去。设 \(p\) 为素数，取 \(R=\mathbb Z_{(p)}\)、\(M=\mathbb Q\)，则 \(p^nM=M\) 对所有 \(n\) 成立，交并不为零。Jacobson 根条件也不能对一般环省略：若 \(k\) 为域，\(k\times k\) 的幂等理想 \(k\times0\) 的各次幂都相同。

\ProseParagraphBreak
本节还给出一种实用判断：设 \((R,\mathfrak m)\) 为 Noether 局部环，\(M\) 为有限生成 \(R\) 模，\(N\subseteq M\) 为子模，\(x\in M\)。若对每个 \(n\ge0\) 都存在 \(y_n\in N\)，使 \(x\equiv y_n\pmod{\mathfrak m^nM}\)，那么 \(x\in N\)，因为 \(N\) 闭。

% !TeX root = ../../Book3.tex
% 纲要标题：### 13.3 完备化
% 纲要位置：工作记录/计划纲要.md，第 2350 行。
% 纲要细目（写作范围）：
% * 逆极限定义
% * 形式幂级数环
% * Noether 情形下有限生成模的完备化正合性
% * 商、有限生成模与张量表示
% * 完备环的 Noether 性；初始形式的有限生成与系数部分和的收敛消去
% * 在完备化局部性质之前证明所需的 Noether 性、完备性与有限生成模的张量表示；把任意逆极限的左正合性与本章有限生成模的完备化正合性区分开

\section{完备化}
\label{b3:sec:1303}


只给出每个精度下的一个答案还不够：高精度答案必须在取商后等于先前的低精度答案。完备化把这种相容性直接写入定义。它首先是一个由商模组成的逆极限；其正合性与 Noether 性则需要另外证明。

\subsection{完备化的逆极限定义}
\label{b3:subsec:1303:01}

\begin{definition}\label{b3:def:inverse-limit-completion}
设 \(R\) 为交换环，\(M_1\xleftarrow{u_1}M_2\xleftarrow{u_2}\cdots\) 为 \(R\) 模及 \(R\) 线性过渡映射组成的逆系统。其\term[niji-xian]{逆极限}[inverse limit]定义为
\[
\varprojlim_nM_n=
\left\{\,(x_n)\in\prod_{n\ge1}M_n\mid u_n(x_{n+1})=x_n\text{ 对所有 }n\,\right\}.
\]
加法与标量作用逐坐标进行。给定理想 \(I\subseteq R\) 与 \(R\) 模 \(M\)，取 \(M_n=M/I^nM\) 和自然商映射，得到 \(M\) 的 \term[wanbeihua]{完备化}[completion]
\[
\widehat M=\varprojlim_n M/I^nM.
\]
自然映射 \(\iota_M:M\rightarrow\widehat M\) 把 \(x\) 送到 \((x+I^nM)_n\)。若此映射为同构，则称 \(M\) 为 \(I\) 进完备模；此处“完备”包含分离性。
\symidx{\(\varprojlim_nM_n\)：模逆系统的逆极限}
\symidx{\(\widehat R\)、\(\widehat M\)：指定理想进拓扑下的完备化}
\end{definition}

逆极限的泛性质可以直接核对：从模 \(P\) 到各 \(M_n\) 的相容映射 \(f_n\)，唯一合成映射 \(p\mapsto\bigl(f_n(p)\bigr)_n\)。另外，\(\ker\iota_M=\bigcap_nI^nM\)。因此自然映射单射所需的正是分离性。

\ProseParagraphBreak
记 \(K_n=\ker(\widehat M\rightarrow M/I^nM)\)。投影 \(\widehat M\rightarrow M/I^nM\) 满射，因为每个余类都可由 \(\iota_M(x)\) 的第 \(n\) 个坐标实现。因此投影诱导的同构
\[
\widehat M/K_n\xrightarrow{\sim}M/I^nM
\]
与过渡商映射相容。在这些同构下，规范映射
\[
\widehat M\longrightarrow\varprojlim_n\widehat M/K_n,
\qquad z\longmapsto(z+K_n)_n
\]
正是 \(\widehat M=\varprojlim_nM/I^nM\) 的恒等映射，故为同构；又 \(\bigcap_nK_n=0\)。

以 \(K_n\) 为邻域基时，这也等价于每个 Cauchy 序列唯一收敛：在每个商中，Cauchy 序列的坐标最终固定，这些固定坐标相容，因而给出唯一极限。反过来，每个相容余类系统逐阶选取代表元便得到 Cauchy 序列，其极限还原这个系统；极限唯一对应分离性。因此这里的逆极限完备性与序列完备性一致。自然映射 \(\iota_M\) 的像稠密，因为任意第 \(n\) 个坐标都可由 \(M\) 中的元素表示。此处的邻域基由投影核 \(K_n\) 构成。

\ProseParagraphBreak
环 \(\widehat R\) 的加法、乘法逐商定义，\(\widehat M\) 是 \(\widehat R\) 模，因为 \((a_n)(m_n)=(a_nm_n)\) 保持相容性。模同态按各阶取商得到完备化同态，并保持恒等映射与复合。

\subsection{形式幂级数环}
\label{b3:subsec:1303:02}

\begin{definition}\label{b3:def:formal-power-series}
设 \(A\) 为交换环，\(r\ge0\) 为整数，\(x_1,\ldots,x_r\) 为形式变量。对 \(\alpha=(\alpha_1,\ldots,\alpha_r)\in\mathbb N^r\)，记 \(x^\alpha=x_1^{\alpha_1}\cdots x_r^{\alpha_r}\)。给定系数族 \((a_\alpha)_{\alpha\in\mathbb N^r}\)，其中 \(a_\alpha\in A\)，把它记为形式和
\(\sum_{\alpha\in\mathbb N^r}a_\alpha x^\alpha\)。全体这样的形式和逐系数相加；对另一个系数族 \((b_\beta)_{\beta\in\mathbb N^r}\)、\(b_\beta\in A\)，按下式相乘：
\[
\left(\sum_\alpha a_\alpha x^\alpha\right)
\left(\sum_\beta b_\beta x^\beta\right)
=\sum_\gamma\left(\sum_{\alpha+\beta=\gamma}a_\alpha b_\beta\right)x^\gamma.
\]
对固定的多重指标 \(\gamma\)，内层和只有有限项，故不需要分析意义的收敛。所得环称为\term[xingshi-miji-shu-huan]{形式幂级数环}[formal power series ring]，记为 \(A[\![x_1,\ldots,x_r]\!]\)。
\end{definition}

交换律、结合律及分配律可在每个固定系数上归结为有限和的同一运算，常数级数 \(1\) 为单位。全次数小于 \(n\) 的截断给出同构
\begin{equation}\label{b3:eq:power-series-limit}
A[\![x_1,\ldots,x_r]\!]\cong
\varprojlim_n A[x_1,\ldots,x_r]/(x_1,\ldots,x_r)^n.
\end{equation}
相容截断唯一指定每个单项式的系数，反过来每个系数族给出相容截断，因此此式既证明双射，也说明环运算为何相同。

\begin{proposition}\label{b3:prop:formal-series-unit}
设 \(A\) 为交换环，\(r\ge0\) 为整数，\(f=\sum_{\alpha\in\mathbb N^r}a_\alpha x^\alpha\in A[\![x_1,\ldots,x_r]\!]\)。则 \(f\) 可逆，当且仅当其常数项 \(a_0\) 在 \(A\) 中可逆。
\end{proposition}
\begin{proof}
若有逆，比较常数项得到必要性。反向写 \(f=a_0(1-h)\)，其中 \(h\) 常数项为零。级数 \(\sum_{j\ge0}h^j\) 有意义，因为全次数小于 \(n\) 的项只受前 \(n\) 项影响。逐阶使用有限等比恒等式得到 \((1-h)\sum_{j\ge0}h^j=1\)，所以 \(f^{-1}=a_0^{-1}\sum_{j\ge0}h^j\)。
\end{proof}

特别地，若 \(k\) 为域，则 \(k[\![x_1,\ldots,x_r]\!]\) 是极大理想为 \((x_1,\ldots,x_r)\) 的局部环。局部多项式环 \(k[x_1,\ldots,x_r]_{(x_1,\ldots,x_r)}\) 也有这个完备化：在每个截断商中，所有原点处不为零的多项式已经可逆，所以先局部化不改变各阶商。若 \(p\) 为素数，则整数环关于 \((p)\) 的完备化为 \(\mathbb Z_p=\varprojlim_n\mathbb Z/p^n\mathbb Z\)。它的进位规则来自这些商环的运算，特征为零，与特征为 \(p\) 的 \(\mathbb F_p[\![t]\!]\) 不同。

\subsection{Noether 环上有限生成模的完备化正合性}
\label{b3:subsec:1303:03}

\begin{lemma}\label{b3:lem:inverse-limit-exact-surjective-kernels}
设 \(R\) 为交换环，给定三个 \(R\) 模逆系统及逐阶短正合列
\(0\rightarrow A_n\rightarrow B_n\rightarrow C_n\rightarrow0\)，其中所有过渡映射与列中的映射交换。若过渡映射 \(A_{n+1}\rightarrow A_n\) 全部满射，则
\[
0\longrightarrow\varprojlim A_n\longrightarrow\varprojlim B_n
\longrightarrow\varprojlim C_n\longrightarrow0
\]
正合。即使不假设满射，左端及中间位置仍正合。
\end{lemma}
\begin{proof}
单射与中间的核可逐坐标检验。要证明右端满射，给定相容的 \((c_n)\)，先选提升 \(b_1\)。若已选 \(b_n\)，任选 \(c_{n+1}\) 的提升 \(b'_{n+1}\)。它投到 \(B_n\) 后与 \(b_n\) 之差落在 \(A_n\)，由满射性可提升成 \(a_{n+1}\in A_{n+1}\)。令 \(b_{n+1}=b'_{n+1}-a_{n+1}\)，就同时保持提升与相容性。递归构造出所需的 \((b_n)\)。
\end{proof}

\begin{theorem}\label{b3:thm:completion-exact}
设 \(R\) 为交换 Noether 环、\(I\subseteq R\) 为理想。对有限生成 \(R\) 模的短正合列
\(0\rightarrow N\rightarrow M\rightarrow P\rightarrow0\)，其 \(I\) 进完备化
\[
0\longrightarrow\widehat N\longrightarrow\widehat M
\longrightarrow\widehat P\longrightarrow0
\]
仍正合。
\end{theorem}
\begin{proof}
把 \(N\) 看作 \(M\) 的子模。逐阶真正正合的列是
\[
0\longrightarrow N/(N\cap I^nM)\longrightarrow M/I^nM
\longrightarrow P/I^nP\longrightarrow0.
\]
其左侧过渡映射为自然商映射，因而满射。使用前一引理取逆极限，右端得到 \(\widehat P\)，中间得到 \(\widehat M\)。左端来自诱导滤过；\cref{b3:thm:artin-rees}给出的
\(I^nN\subseteq N\cap I^nM\subseteq I^{n-c}N\)
说明该滤过与 \(N\) 自身的 \(I\) 幂滤过相互共尾。由\cref{b3:prop:stable-filtration-equivalent}的自然同构，左端识别为 \(\widehat N\)，得到结论。
\end{proof}

这里没有断言任意逆极限右正合。例如短正合列
\(0\rightarrow2^n\mathbb Z\rightarrow\mathbb Z\rightarrow\mathbb Z/2^n\mathbb Z\rightarrow0\)
的逆极限右端为 \(\mathbb Z_2\)，但 \(\mathbb Z\rightarrow\mathbb Z_2\) 不满：每个模 \(2^n\) 的环中 \(3\) 的逆元组成相容系统，若它来自整数 \(a\)，则整数 \(3a-1\) 被所有 \(2^n\) 整除，从而等于零，矛盾。本例的左侧过渡映射为严格包含，不满足引理的满射假设。

\subsection{商、有限生成模与张量表示}
\label{b3:subsec:1303:04}

以下 \(R\) 均为 Noether 环，\(I\) 为固定理想，\(M\) 为有限生成模。

\begin{proposition}\label{b3:prop:completion-quotients}
设 \(R\) 为交换 Noether 环、\(I,J\subseteq R\) 为理想、\(M\) 为有限生成 \(R\) 模，所有完备化均关于 \(I\) 进行。在 \(\widehat M\) 中有
\[
\operatorname{im}(\widehat{JM}\rightarrow\widehat M)=J\widehat M,
\qquad \widehat{M/JM}\cong\widehat M/J\widehat M.
\]
特别地，
\[
\widehat M/I^n\widehat M\cong M/I^nM,
\qquad \ker(\widehat M\rightarrow M/I^nM)=I^n\widehat M.
\]
因此 \(\widehat M\) 对 \(I\widehat R\) 进拓扑完备且分离。
\end{proposition}
\begin{proof}
取 \(J=(a_1,\ldots,a_s)\)。映射
\(M^s\rightarrow JM\)、\((m_j)\mapsto\sum_ja_jm_j\) 满射。由完备化正合性，其完备化仍满；有限直和的完备化逐坐标等于完备化的直和，所以它在 \(\widehat M\) 中的像恰为 \(\sum_ja_j\widehat M=J\widehat M\)。再完备化 \(0\rightarrow JM\rightarrow M\rightarrow M/JM\rightarrow0\)，得到商的同构。

取 \(J=I^n\)。模 \(M/I^nM\) 在第 \(n\) 阶以后各商都不再改变，故其完备化是自身。于是得到第二组公式；逆极限原有的邻域核正是理想的幂，前面已证的完备性和分离性遂成为理想进完备性与分离性。
\end{proof}

\begin{theorem}\label{b3:thm:completion-tensor}
设 \(R\) 为交换 Noether 环、\(I\subseteq R\) 为理想、\(M\) 为有限生成 \(R\) 模，\(\widehat R\) 与 \(\widehat M\) 均表示 \(I\) 进完备化。则自然映射
\[
\widehat R\otimes_RM\longrightarrow\widehat M,
\qquad a\otimes m\longmapsto a\,\iota_M(m)
\]
是同构。特别地，\(\widehat M\) 为有限生成 \(\widehat R\) 模。
\end{theorem}
\begin{proof}
\(R\) 为 Noether 环，故 \(M\) 有有限展示
\(R^u\xrightarrow{A}R^v\rightarrow M\rightarrow0\)。完备化正合性说明 \(\widehat M\) 为矩阵 \(A:\widehat R^u\rightarrow\widehat R^v\) 的余核。张量积的右正合性说明 \(\widehat R\otimes_RM\) 为同一矩阵的余核。自然映射在自由模处就是逐坐标同构，故在余核处也是同构；这里比较余核只需张量积的右正合性。
\end{proof}

\subsection{完备环的 Noether 性}
\label{b3:subsec:1303:noetherian}

商的同构还保留相邻两阶之间的关系，从而识别原环与完备环的伴随分次环。分次环中有限个初始形式能够控制一个理想的全部初始形式；要回到原理想，则需把逐阶消去产生的系数收敛到完备环中。

\begin{theorem}\label{b3:thm:completion-noetherian}
设 \(R\) 为交换 Noether 环、\(I\subseteq R\) 为理想，\(\widehat R\) 为 \(I\) 进完备化。则 \(\widehat R\) 是 Noether 环，并且
\[
\operatorname{gr}_{I\widehat R}(\widehat R)
\cong\operatorname{gr}_I(R).
\]
\end{theorem}
\begin{proof}
记 \(B=\widehat R\)、\(K=IB\)。由\cref{b3:prop:completion-quotients}，\(B\) 对 \(K\) 进拓扑完备且分离。该命题的商同构在相邻两阶之间相容，取核得到 \(K^n/K^{n+1}\cong I^n/I^{n+1}\)，且乘法相容。若 \(I=(a_1,\ldots,a_s)\)，此分次环是 \((R/I)[X_1,\ldots,X_s]\) 的商，故为 Noether 环。

设 \(J\subseteq B\) 为任意理想。对每个非零 \(f\in B\)，分离性使它有唯一阶数 \(d\ge0\)，即 \(f\in K^d\setminus K^{d+1}\)；记其初始形式为 \(\operatorname{in}(f)\)。由所有 \(f\in J\) 的初始形式生成的齐次理想有限生成，故可选有限个 \(f_1,\ldots,f_r\in J\)，使它们的初始形式生成该理想，记阶数为 \(d_i\)。这一步可以选实际初始形式，是因为原生成族中有限个已足以生成同一个理想。

对 \(0\ne f\in J\)，令 \(r_0=f\)。若第 \(\nu\) 次的余项 \(r_\nu\ne0\)，记其阶数为 \(n_\nu\)，将初始形式表示成 \(\sum_i\overline a_{i,\nu}\operatorname{in}(f_i)\)。当 \(n_\nu\ge d_i\) 时，取 \(\overline a_{i,\nu}\) 为次数 \(n_\nu-d_i\) 的齐次元，并提升为 \(a_{i,\nu}\in K^{n_\nu-d_i}\)；当 \(n_\nu<d_i\) 时，取 \(a_{i,\nu}=0\)。于是
\[
r_{\nu+1}=r_\nu-\sum_i a_{i,\nu}f_i
\in J\cap K^{n_\nu+1}.
\]
若余项为零，已经得到有限组合；否则 \(n_{\nu+1}>n_\nu\)。若此过程无限继续，则 \(n_\nu\to\infty\)。设
\[
b_{i,\nu}=\sum_{\mu=0}^{\nu-1}a_{i,\mu},
\qquad r_\nu=f-\sum_i b_{i,\nu}f_i.
\]
对任意 \(N\ge0\)，当 \(\nu\) 充分大时，\(n_\nu-d_i\ge N\) 对每个 \(i\) 成立，所以以后的全部系数增量均属于 \(K^N\)。因此 \((b_{i,\nu})_\nu\) 是 Cauchy 序列，由 \(B\) 的完备性收敛于某个 \(b_i\in B\)。同时 \(r_\nu\in K^{n_\nu}\) 趋于零。对每个 \(N\)，在充分大的 \(\nu\) 处有 \(b_i-b_{i,\nu}\in K^N\) 及 \(r_\nu\in K^N\)，故
\[
f-\sum_i b_if_i
=r_\nu+\sum_i(b_{i,\nu}-b_i)f_i\in K^N.
\]
由 \(\bigcap_NK^N=0\) 得 \(f=\sum_i b_if_i\)。因此 \(J=(f_1,\ldots,f_r)\)；零理想同样有限生成，故 \(B\) 为 Noether 环。
\end{proof}

这个证明同时说明为什么“完备”有实质作用：初始形式只能逐阶消去，系数最后可能成为无穷和。有限阶消去若没有收敛到环中元素的保证，就不能推出原理想由所选元素生成。

% !TeX root = ../../Book3.tex
% 纲要标题：### 13.4 完备化的局部性质
% 纲要位置：工作记录/计划纲要.md，第 2381 行。
% 纲要细目（写作范围）：
% * 完备化与张量积
% * Noether 局部环完备化的忠实平坦性
% * 极大理想、剩余域及维数
% * 用依赖图连接 Rees 有限性、Artin–Rees、完备化正合性、张量表示与平坦性；以有限阶商的长度比较说明维数与重数保持
% * 证明顺序固定为：有限生成模的完备化正合性、\(M\otimes_R\widehat R\cong\widehat M\)、第12.1节的理想张量与局部判据、平坦性、剩余域非零及第12.5节的忠实平坦性判据
% * 极大理想进完备化的维数比较使用第11章的伴随分次与 Hilbert–Samuel 理论；模的忠实平坦下降仍是独立选读，不作为此处证明前提

\section{完备化的局部性质}
\label{b3:sec:1304}


完备化把原环中的计算搬到一个容纳所有相容近似的环中。上一节的正合性与张量表示使我们能够比较两边的模同态：有限生成理想的包含在完备化后仍为单射，因而完备环对原环平坦。在局部环中，剩余域还保持不变；这一点使变基后的计算能够反向检测原来的模。

\subsection{完备化与张量积}
\label{b3:subsec:1304:01}

设 \(R\) 为交换 Noether 环，\(I\subseteq R\) 为理想，各完备化均关于 \(I\) 进行。对有限生成 \(R\) 模 \(M,N\)，\cref{b3:thm:completion-tensor}给出的同构与模同态相容：若 \(f:M\rightarrow N\)，两边都把 \(a\otimes m\) 送到 \(a\,\iota_N\bigl(f(m)\bigr)\)。因此可以把有限生成模之间的完备化映射统一视为沿 \(R\rightarrow\widehat R\) 的标量扩张。

\begin{proposition}\label{b3:prop:completion-finite-tensor-products}
设 \(R\) 为交换 Noether 环、\(I\subseteq R\) 为理想，\(M,N\) 为有限生成 \(R\) 模；记 \(\widehat R\)、\(\widehat M\)、\(\widehat N\) 及 \(\widehat{M\otimes_RN}\) 为各自的 \(I\) 进完备化。则
\[
\widehat{M\otimes_RN}
\cong\widehat M\otimes_{\widehat R}\widehat N.
\]
\end{proposition}
\begin{proof}
左边由张量表示等于 \(\widehat R\otimes_R(M\otimes_RN)\)。右边把两因子写成 \(\widehat R\otimes_RM\)、\(\widehat R\otimes_RN\)，映射
\[
(a\otimes m)\otimes(b\otimes n)\longmapsto ab\otimes(m\otimes n)
\]
保持各个平衡关系，逆映射为
\(c\otimes(m\otimes n)\mapsto(c\otimes m)\otimes(1\otimes n)\)。在纯张量上检查两复合为恒等，便得同构。
\end{proof}

对于不有限生成的模，这个张量表示可能失效。例如设 \(p\) 为素数，取 \(R=\mathbb Z_{(p)}\)、\(I=(p)\)、\(M=\mathbb Q\)。有限阶商 \(R/p^nR\cong\mathbb Z/p^n\mathbb Z\) 给出 \(\widehat R=\mathbb Z_p\)。由于 \(p^nM=M\)，有 \(\widehat M=0\)；但
\[
\widehat R\otimes_RM=\mathbb Z_p[1/p]\ne0.
\]
由局部化的零判据，\(\mathbb Z_p[1/p]=0\) 当且仅当某个 \(p^n\) 在 \(\mathbb Z_p\) 中为零。但对每个 \(n\ge0\)，\(p^n\) 在第 \(n+1\) 阶坐标 \(\mathbb Z/p^{n+1}\mathbb Z\) 中非零，故在逆极限中也非零。因此倒置 \(p\) 后仍为非零环，张量积与完备化在这个模上不同。

\subsection{Noether 局部环完备化的忠实平坦性}
\label{b3:subsec:1304:02}
\label{b3:subsec:1304:04}

\begin{theorem}\label{b3:thm:completion-flat}
若 \(R\) 为交换 Noether 环、\(I\subseteq R\) 为理想，则其 \(I\) 进完备化 \(\widehat R\) 是平坦 \(R\) 模。若 \(R\) 还是局部环，极大理想为 \(\mathfrak m\)，且 \(I\subseteq\mathfrak m\)，则 \(R\rightarrow\widehat R\) 忠实平坦。特别地，极大理想进完备化忠实平坦。
\end{theorem}
\begin{proof}
对任意有限生成理想 \(J\subseteq R\)，映射
\[
J\otimes_R\widehat R\longrightarrow R\otimes_R\widehat R=\widehat R
\]
由自然张量表示等同于 \(\widehat J\rightarrow\widehat R\)。\cref{b3:thm:completion-exact}说明后一个映射单射。于是\cref{b3:thm:ideal-flatness-criterion}给出平坦性。

若 \(R\) 局部且 \(I\subseteq\mathfrak m\)，则 \(R/\mathfrak m\) 的 \(I\) 进滤过第一步就为零。由\cref{b3:prop:completion-quotients}，
\[
\widehat R/\mathfrak m\widehat R\cong R/\mathfrak m\ne0.
\]
这说明完备化没有删去原谱的唯一闭点。\(R\) 只有这一个极大理想，故结合已经证明的平坦性，\cref{b3:thm:faithfully-flat-criterion}给出忠实平坦性，并保证所有素点都有上方点。
\end{proof}

\cref{b3:fig:completion-flat-dependencies}把上述论证所用的结果连在一起。图中“正合”限于有限生成模的完备化；得到环同态的平坦性后，张量积保持正合的结论才适用于任意模。

\begin{center}
\begin{minipage}{\textwidth}
\makeatletter
\def\@captype{figure}
\makeatother
\centering
\[
\begin{tikzcd}[column sep=large,row sep=large]
\boxed{\mathcal R_I(R)\;\text{是 Noether 环}}
  \arrow[r]
&\boxed{\text{Artin--Rees}}
  \arrow[r]
&\boxed{\text{有限生成模的完备化正合}}
  \arrow[d]\\
\boxed{R\rightarrow\widehat R\;\text{忠实平坦}}
&\boxed{R\rightarrow\widehat R\;\text{平坦}}
  \arrow[l,"R\;\text{局部，}I\subseteq\mathfrak m"']
&\boxed{\widehat M\cong\widehat R\otimes_RM}
  \arrow[l,"\text{理想张量判据}"']
\end{tikzcd}
\]
\caption[完备化的正合性、张量表示与平坦性]{完备化的正合性、张量表示与平坦性。设 \(R\) 为交换 Noether 环，\(I\subseteq R\) 为理想，\(M\) 为有限生成 \(R\) 模；局部情形以 \(\mathfrak m\) 表示极大理想。}
\label{b3:fig:completion-flat-dependencies}
\end{minipage}
\end{center}

在 \(R\) 局部且 \(I\subseteq\mathfrak m\) 的情形中，忠实平坦性意味着：一个 \(R\) 模即使不是有限生成的，也不会在张量到 \(\widehat R\) 后从非零变成零；模映射的单射、满射及正合性也能在张量后检测。这里谈的是张量积函子。上一小节的 \(M=\mathbb Q\) 例子表明，对任意模不能把它替换成完备化函子。

\ProseParagraphBreak
全局完备化未必忠实。例如设 \(p,q\) 为不同素数，\(\mathbb Z\rightarrow\mathbb Z_p\) 平坦，但 \(q\) 在 \(\mathbb Z_p\) 中可逆：在各个 \(\mathbb Z/p^n\mathbb Z\) 中，\(q\) 的唯一逆元组成相容系统。因此
\[
(\mathbb Z/q\mathbb Z)\otimes_{\mathbb Z}\mathbb Z_p=0.
\]
几何上，\((q)\) 没有上方素点，因为 \(q\) 可逆，不能属于 \(\mathbb Z_p\) 的任何真理想。先局部化为 \(\mathbb Z_{(p)}\) 再完备化，则由前述定理得到忠实平坦映射；其谱映射满射，而 \(\operatorname{Spec}\mathbb Z_{(p)}\) 在 \(\operatorname{Spec}\mathbb Z\) 中对应 \((0),(p)\)。因此
\[
\operatorname{im}\bigl(\operatorname{Spec}\mathbb Z_p\to\operatorname{Spec}\mathbb Z\bigr)
=\{(0),(p)\}.
\]
关于合数 \(m\) 的 \((m)\) 进完备化，\cref{b3:ex:13-composite-adic}要求通过各有限阶商的中国剩余分解比较不同素数的分量。局部性与忠实性在这种全局完备化中仍须分别检查。

\subsection{极大理想、剩余域与维数}
\label{b3:subsec:1304:03}

\begin{proposition}\label{b3:prop:completed-local-ring}
设 \((R,\mathfrak m)\) 为交换 Noether 局部环，\(k=R/\mathfrak m\) 为剩余域，\(B=\widehat R\) 为极大理想进完备化。则 \(B\) 是 Noether 局部环，其极大理想为 \(\mathfrak n=\mathfrak mB\)，剩余域自然等于 \(k\)。对每个 \(r\ge1\) 有
\[
B/\mathfrak n^r\cong R/\mathfrak m^r,
\qquad \operatorname{gr}_{\mathfrak n}(B)\cong\operatorname{gr}_{\mathfrak m}(R).
\]
\end{proposition}
\begin{proof}
Noether 性和两项同构分别由\cref{b3:thm:completion-noetherian}和\cref{b3:prop:completion-quotients}给出。商 \(B/\mathfrak n=k\) 是域，所以 \(\mathfrak n\) 为极大理想。若 \(b\notin\mathfrak n\)，选 \(a\in R\) 使 \(ab\equiv1\pmod{\mathfrak n}\)，并写 \(ab=1-u\)、\(u\in\mathfrak n\)。因为 \(u^j\in\mathfrak n^j\)，部分和 \(s_N=1+u+\cdots+u^N\) 在模每个固定 \(\mathfrak n^r\) 后最终不变，所以是 \(\mathfrak n\) 进 Cauchy 序列。\(B\) 的完备性给出极限 \(s=\sum_{j\ge0}u^j\)；有限等比恒等式 \((1-u)s_N=1-u^{N+1}\) 在每个有限阶商中给出 \((1-u)s=1\)。于是 \(b(as)=1\)，故 \(b\) 可逆。这与\cref{b3:prop:formal-series-unit}中的逆元构造使用同一个逐阶计算。因此 \(\mathfrak n\) 之外的每个元素均为单位；而 \(\mathfrak n\) 是真理想，其中不含单位，所以它恰由所有非单位组成，是 \(B\) 的唯一极大理想。
\end{proof}

所有有限阶商及剩余域保持不变，并不表示原环与完备环相同。例如设 \(k\) 为域，\(R=k[x]_{(x)}\) 的元素是原点分母非零的有理函数，而 \(\widehat R=k[\![x]\!]\) 中的级数 \(v=\sum_{j\ge0}x^{2^j}\) 不在 \(R\) 中。若 \(v=P/Q\)，其中 \(P,Q\in k[x]\)、\(Q(0)\ne0\)，则对充分大的 \(j\)，\(2^j>\deg P\) 且 \(2^{j-1}>\deg Q\)。在 \(Qv=P\) 中比较 \(x^{2^j}\) 的系数，左侧仅有 \(Q(0)\) 的贡献，右侧为零，矛盾。尽管对每个 \(n\) 都有 \(R/(x^n)\cong\widehat R/(x^n)\)，这个新元素仍需要一整套无限相容的截断才能指定；任一固定阶都可以用原环的多项式表示它。维数保持则由下一小节的长度增长比较给出。

\subsection{Hilbert–Samuel 理论与完备化的维数保持}
\label{b3:subsec:1304:05}

\begin{theorem}\label{b3:thm:completion-dimension}
对交换 Noether 局部环 \((R,\mathfrak m)\) 及其极大理想进完备化 \(B=\widehat R\)、\(\mathfrak n=\mathfrak mB\)，有
\[
\dim B=\dim R,\qquad e_{\mathfrak n}(B)=e_{\mathfrak m}(R).
\]
更一般地，对非零有限生成的 \(R\) 模 \(M\)，
\[
\dim_B\widehat M=\dim_RM,\qquad
 e_{\mathfrak n}(\widehat M)=e_{\mathfrak m}(M).
\]
\end{theorem}
\begin{proof}
\cref{b3:prop:completion-quotients}给出
\(\widehat M/\mathfrak n^{r+1}\widehat M\cong M/\mathfrak m^{r+1}M\)。对每个整数 \(r\ge0\)，两者的标量作用都通过同一个商环 \(B/\mathfrak n^{r+1}\cong R/\mathfrak m^{r+1}\) 分解，所以子模及组成列完全相同，特别是
\[
\ell_B(\widehat M/\mathfrak n^{r+1}\widehat M)
=\ell_R(M/\mathfrak m^{r+1}M).
\]
因此两侧的 Hilbert–Samuel 函数逐项相同。\(B\) 已证为 Noether 局部环，\(\widehat M\) 已证为有限生成 \(B\) 模且由忠实平坦性非零，故两侧都能应用\cref{b3:thm:hilbert-samuel}。相等的累计函数给出相等的最终多项式，比较次数与首项即得维数和重数。取 \(M=R\) 得环的结论。
\end{proof}

因此完备化保留局部长度增长及其决定的维数，却未必保留素理想分解的所有细节。例如曲线在形式邻域中可能分出原环中尚未分开的分支；维数相等并不声称两者素谱同胚。

% !TeX root = ../../Book3.tex
% 纲要标题：### 13.5 有限阶近似与形式邻域
% 纲要位置：工作记录/计划纲要.md，第 2390 行。
% 纲要细目（写作范围）：
% * 商 \(R/\mathfrak m^N\)、截断数据及相容逆系统
% * 有限阶计算、完整形式对象与收敛解析对象的不同信息量
% * Noether 局部环中理想向完备化扩张再收缩的保持；与忠实平坦性的联系
% * 由子模闭性证明线性方程可解性从完备化下降，并以非线性形式根不能下降的例子说明适用范围；习题比较解集、精确解的有限阶逼近与分支数据
% * 算法须注明多项式、代数幂级数或有限截断的输入形式；不能对任意不可计算幂级数声称有限有效求解
% ---

\section{有限阶近似与形式邻域}
\label{b3:sec:1305}


完备化由全部有限阶商组成，但一次计算通常只能得到其中有限多个。本节说明这类计算究竟确定什么，以及何时能够从所有阶的近似恢复原环中的结论。

\subsection{有限阶商、截断数据与相容逆系统}
\label{b3:subsec:1305:01}

\begin{definition}\label{b3:def:finite-order-neighbourhood}
设 \((R,\mathfrak m)\) 为交换 Noether 局部环，\(N\ge0\) 为整数。本书将商环 \(R/\mathfrak m^{N+1}\) 称为该点的 \(N\) 阶邻域的坐标环。
\end{definition}

指数取 \(N+1\)，是为了保留直到第 \(N\) 阶的信息：零阶商 \(R/\mathfrak m\) 只剩剩余域；一阶商 \(R/\mathfrak m^2\) 还保留 \(\mathfrak m/\mathfrak m^2\) 中的一次信息；二阶商再保留 \(\mathfrak m^2/\mathfrak m^3\) 中的二次信息。所有商 \(R/\mathfrak m^q\)、\(q\ge1\)，连同它们之间的自然商映射组成相容逆系统，描述此点的\term[xingshi-linyu]{形式邻域}[formal neighbourhood]。

\ProseParagraphBreak
设 \(k\) 为域。在 \(k[x,y]_{(x,y)}\) 中模掉 \((x,y)^3\)，得到 \(k\) 基
\(1,x,y,x^2,xy,y^2\)。相乘后删去全次数至少三的单项式，即为商环乘法。对尖点局部环
\(R=k[x,y]_{(x,y)}/(y^2-x^3)\)、\(\mathfrak m=(x,y)R\)，二阶邻域的坐标环为
\[
R/\mathfrak m^3\cong k[x,y]/\bigl((x,y)^3,y^2\bigr).
\]
最低关系 \(y^2=0\) 与\cref{b3:prop:hypersurface-tangent-cone}所给切锥的关系 \(Y^2=0\) 相同，但这个有限阶商还截去了全次数至少三的全部单项式。在三阶邻域 \(R/\mathfrak m^4\) 中，右侧的三次项 \(x^3\) 才重新可见。一个低阶商无法单独区分所有高阶方程。

\begin{proposition}\label{b3:prop:compatible-formal-solutions}
设 \((R,\mathfrak m)\) 为交换 Noether 局部环，\(\widehat R\) 为极大理想进完备化，\(f_1,\ldots,f_s\in R[X_1,\ldots,X_r]\)。在 \(\widehat R^r\) 中给定一个共同零点，等价于给定对每个 \(N\ge1\) 的零点
\(a_N\in(R/\mathfrak m^N)^r\)，且 \(a_{N+1}\) 模 \(\mathfrak m^N\) 等于 \(a_N\)。
\end{proposition}
\begin{proof}
零点逐阶取商显然相容。反向由逆极限定义，各坐标的相容余类给出 \(a\in\widehat R^r\)。多项式运算与取商交换，所以每个 \(f_i(a)\) 在所有有限阶商中为零；\(\widehat R\) 分离，故 \(f_i(a)=0\)。
\end{proof}

逐阶任意选取零点 \(a_N\)，只保证每个有限阶的方程成立，未必保证这些选择彼此相容。形式解要求一次选出整个系统，并同时满足
\[
\exists(a_N)_{N\ge1}\quad
\begin{cases}
f_i(a_N)=0\text{ 于 }R/\mathfrak m^N
 &\text{对所有 }i,N,\\
a_{N+1}\bmod\mathfrak m^N=a_N
 &\text{对所有 }N.
\end{cases}
\]
是否能够作出这样的相容选择，还须利用环与方程的具体条件。这与\secref{b3:sec:1303}中逐阶满射未必给出逆极限满射的区别相联系；简单根的逐阶提升将在\chapref{b3:ch:23}系统讨论。

\subsection{有限阶、形式与收敛解析对象}
\label{b3:subsec:1305:02}

在复数域上，多项式、原点附近收敛的幂级数与任意形式幂级数分别给出
\[
\mathbb C[x]\subseteq\mathbb C\{x\}\subseteq\mathbb C[\![x]\!].
\]
这里多项式只允许有限多个非零系数，\(\mathbb C\{x\}\) 允许无限多个系数，但要求正收敛半径；形式幂级数则允许任意系数序列，乘法只作逐阶有限的系数运算。\(\sum_{n\ge0}x^n\) 在 \(|x|<1\) 时收敛，却不是多项式；\(\sum_{n\ge0}n!x^n\) 是形式级数。对任意固定的 \(x\in\mathbb C\setminus\{0\}\)，令 \(a_n=n!x^n\)，则
\[
\dfrac{|a_{n+1}|}{|a_n|}=(n+1)|x|\longrightarrow\infty.
\]
因此从某项起 \(|a_{n+1}|\ge2|a_n|\)，通项 \(n!x^n\) 不趋于零，级数不收敛。

\ProseParagraphBreak
任意给定的形式级数模 \(x^N\) 都由一个次数小于 \(N\) 的多项式实现。相同的截断既可接上全零的尾部，成为一个收敛多项式，也可接上从第 \(N\) 项开始的阶乘系数尾部，成为发散的形式级数。因此有限阶数据不能判定收敛性。完备化是代数上的逐阶补全，不自动赋予解析意义；若研究收敛性，还必须另加系数估计或相应的解析定理。

\subsection{完备化中理想的扩张与收缩}
\label{b3:subsec:1305:03}

\begin{proposition}\label{b3:prop:completion-ideal-contraction}
设 \((R,\mathfrak m)\) 为交换 Noether 局部环，\(B=\widehat R\) 为极大理想进完备化。对每个理想 \(J\subseteq R\)，有
\[
JB\cap R=\bigcap_{n\ge1}(J+\mathfrak m^n)=J,
\qquad B/JB\cong\widehat{R/J}.
\]
后一完备化关于 \((\mathfrak m+J)/J\) 进行。更一般地，设 \(M\) 为有限生成 \(R\) 模，\(N\subseteq M\) 为子模，各模均取 \(\mathfrak m\) 进完备化。自然映射 \(M\to\widehat M\) 及 \(\widehat N\to\widehat M\) 均单射；经这些映射将它们识别为子模后，有
\[
M\cap\widehat N=N\quad\text{于 }\widehat M.
\]
因此，对任意非负整数 \(r,s\)、\(R\) 上的 \(s\times r\) 矩阵 \(A\) 及 \(b\in R^s\)，线性方程 \(Ax=b\) 在 \(B^r\) 中有解，当且仅当在 \(R^r\) 中有解。

一般非线性方程没有这一可解性下降结论。若 \(k\) 为特征不为二的域，\(x\) 为不定元，则 \(k[\![x]\!]\) 中存在唯一常数项为一、满足 \(u^2=1+x\) 的级数；但在 \(k=\mathbb Q\) 时，此方程在原局部环 \(\mathbb Q[x]_{(x)}\) 中无解。
\end{proposition}
\begin{proof}
\(\mathfrak m\) 是 \(R\) 的 Jacobson 根，故\cref{b3:thm:krull-intersection}使每个有限生成模到其完备化的自然映射单射。Noether 性使 \(N\) 与 \(M/N\) 均有限生成；对 \(0\rightarrow N\rightarrow M\rightarrow M/N\rightarrow0\) 应用\cref{b3:thm:completion-exact}，其完备化仍正合。因此 \(\widehat N\) 可识别为 \(\widehat M\) 的子模。一个 \(x\in M\) 的像属于 \(\widehat N\)，当且仅当它在 \(\widehat{M/N}\) 中为零；由 \(M/N\to\widehat{M/N}\) 的单射性，这又等价于 \(x\in N\)，即得交式。

取 \(M=R\)、\(N=J\)，\cref{b3:prop:completion-quotients}将 \(\widehat J\) 的像识别为 \(JB\)，并给出 \(B/JB\cong\widehat{R/J}\)，故得到理想的扩张收缩公式。同一个 Krull 交结论对有限生成 \(R\) 模 \(R/J\) 给出 \(\bigcap_{n\ge1}\mathfrak m^n(R/J)=0\)；取逆像就是 \(\bigcap_{n\ge1}(J+\mathfrak m^n)=J\)。于是，只要原环元素在每个有限精度下都能被 \(J\) 中元素逼近，它本来就属于 \(J\)。

对线性方程，令 \(M=R^s\)、\(N=\operatorname{im}A\)。\(N\) 由 \(A\) 的有限列生成，满射 \(R^r\to N\) 的核也由 Noether 性有限生成。对这个满射及包含 \(N\to R^s\) 分别使用完备化正合性，得到 \(\operatorname{im}(\widehat A)=\widehat N\) 在 \(B^s\) 中的像。因此完备环中有解，恰说明 \(b\) 的像属于 \(\widehat N\)；刚证的交式使它等价于 \(b\in N\)，也就是原环中有解。这是解的存在性结论，选定的形式解仍可在 \(\ker\widehat A\) 中任意改变，并不必来自 \(R^r\)；例如零方程 \(0x=0\) 的每个形式级数都是解。

对于非线性反例，写 \(u=1+\sum_{n\ge1}a_nx^n\)。比较 \(u^2=1+x\) 的系数，得到
\[
2a_1=1,\qquad
2a_n=-\sum_{i=1}^{n-1}a_ia_{n-i}\quad(n\ge2).
\]
因 \(2\) 在 \(k\) 中可逆，这个递归唯一确定所有系数，并使每一阶的方程成立，所以给出唯一常数项为一的形式根。若 \(k=\mathbb Q\) 时有原环中的根，将它写成互素非零多项式之比 \(P/Q\)，就有 \(P^2=(1+x)Q^2\)。不可约因子 \(1+x\) 在左侧的指数为偶数，在右侧的指数为奇数，矛盾。因此连 \(\mathbb Q(x)\) 中也没有根，而 \(\mathbb Q[x]_{(x)}\) 的完备化 \(\mathbb Q[\![x]\!]\) 中有上述形式根。
\end{proof}

Krull 交定理把理想收缩解释为所有阶近似所检验的闭性；另一方面，\cref{b3:thm:completion-flat}已证明 \(R\to B\) 忠实平坦，故同一收缩等式也是\cref{b3:prop:faithfully-flat-detection}第二项的一次应用。

\subsection{幂级数计算的输入形式与有效算法的适用范围}
\label{b3:subsec:1305:04}

设 \(k\) 为域，\(r,N\ge1\) 为整数，记 \(A=k[\![x_1,\ldots,x_r]\!]\)、\(\mathfrak n=(x_1,\ldots,x_r)A\)。当输入为 \(k[x_1,\ldots,x_r]\) 中的多项式并指定精度 \(N\) 时，模 \(\mathfrak n^N\) 的加法、乘法及求单位的逆只涉及有限个系数。若 \(k\) 中的系数具有有限表示，且四则运算与相等判定可有效执行，这些是有限算法。例如对常数项为一的 \(f=1-h\)，有 \(h\in\mathfrak n\)，其模 \(\mathfrak n^N\) 的逆可取
\[
1+h+\cdots+h^{N-1}\pmod{\mathfrak n^N},
\]
因为 \((1-h)(1+h+\cdots+h^{N-1})=1-h^N\)，而 \(h^N\in\mathfrak n^N\)。例如\cref{b3:ex:13-finite-inverse}中的 \(1-x-y^2\) 模 \((x,y)^4\) 的逆，就由这个有限和截去全次数至少四的项算出。

\begin{definition}\label{b3:def:algebraic-series-input}
设 \(k\) 为域，\(r\ge1\) 为整数，\(x_1,\ldots,x_r\) 为不定元，\(f\in k[\![x_1,\ldots,x_r]\!]\)。若 \(f\) 在有理函数域 \(k(x_1,\ldots,x_r)\) 上代数，即存在非零多项式 \(P\in k[x_1,\ldots,x_r,T]\) 使 \(P(x_1,\ldots,x_r,f)=0\)，则称 \(f\) 为\term[algebraicpowerseries]{代数幂级数}[algebraic power series]。
\end{definition}

若输入是代数幂级数，必须同时给出定义方程、所选分支及取得所需系数的方法。例如特征不为二时，\cref{b3:prop:completion-ideal-contraction}中的方程 \(T^2-(1+x)=0\) 有两个形式分支 \(u,-u\)，常数项分别为 \(1,-1\)。只给这个方程尚未指定输入的是哪一个级数；指定常数项后，该例的递归才给出所需有限阶系数。一般定义方程还可能有不能用选定参数表示成幂级数的分支，取得系数的方法须针对所选分支说明。任意形式级数则含有无限数据，未必具有有限编码或可计算的系数序列；它作为逆极限中的元素存在，并不自动给出可供有限算法操作的输入。完备性保证相容近似有极限，对象的有限编码与精确判等还需要另外的可计算性条件。

\ProseParagraphBreak
如果获取一个任意形式级数的信息只靠逐个查询系数，有限次查询不能完成对所有输入的零判定。事实上，设程序在零级数上查询的各单项式全次数至多为 \(N\)，则 \(0\) 与 \(x_1^{N+1}\) 在所有已查询位置给出相同答案，程序的这些查询无法区分二者。因此截断计算给出的是指定精度内的确定结论；要从中推断完整等式，还须有额外的次数界、有限生成关系或已经证明的提升与唯一性结果。局部标准基如何提供这类有限控制，将在本卷附录 F 中进一步讨论。


\begin{chaptersummary}
Artin–Rees 的关键在于把 Noether 环上无限多层子模组成一个有限生成的 Rees 模。它给出诱导滤过的稳定性，使子模自身的拓扑与环境诱导的拓扑一致；配合 Nakayama 引理，便得到 Jacobson 根条件下有限生成模的分离性和子模闭性。

\ProseParagraphBreak
完备化的正合性先在逐阶商中建立，再通过相容提升和滤过比较得到。Noether 性使有限生成模具有有限展示，从而可用同一有限矩阵计算完备化与 \(\widehat R\) 标量扩张，得到 \(\widehat M\cong\widehat R\otimes_RM\)。理想张量判据随后给出平坦性，非零剩余域给出局部忠实性。完备环的 Noether 性由初始形式的有限生成与收敛消去证明，而维数保持由相同的 Hilbert–Samuel 函数证明。

\ProseParagraphBreak
形式对象保存所有相容截断；单个截断只保存有限精度。局部理想的闭性允许从任意高精度的理想成员关系恢复原环中的成员关系，却不能据此把任意形式级数视为可计算或收敛的函数。
\end{chaptersummary}

\BookExercises

\begin{exercise}\label{b3:ex:13-equivalent-ideals}
设 \(R\) 为交换 Noether 环，\(I,J\) 为 \(R\) 的理想且满足 \(\sqrt I=\sqrt J\)。证明任意 \(R\) 模上的 \(I\) 进拓扑与 \(J\) 进拓扑相同，且其完备化自然同构。
\end{exercise}
\begin{solution}
取 \(I=(a_1,\ldots,a_r)\)，各 \(a_i\) 的某个幂属于 \(J\)。共同选择整数 \(c\ge1\) 使 \(a_i^c\in J\)，并令 \(a=r(c-1)+1\)。所有全次数至少 \(a\) 的单项式都有一个指数至少 \(c\)，故 \(I^a\subseteq J\)。交换 \(I,J\) 同理得某个整数 \(b\ge1\) 使 \(J^b\subseteq I\)。于是 \(I^{an}M\subseteq J^nM\)、\(J^{bn}M\subseteq I^nM\)，两组邻域相互共尾。将第 \(an\) 阶相容余类投到 \(M/J^nM\)，以及将第 \(bn\) 阶相容余类投到 \(M/I^nM\)，构成互逆的自然极限同构。
\end{solution}

\begin{exercise}\label{b3:ex:13-nonclosed-submodule}
在 \(\mathbb Z\) 的 \((2)\) 进拓扑中求 \(3\mathbb Z\) 的闭包。比较在 \(\mathbb Z_{(2)}\) 中相应理想的情形。
\end{exercise}
\begin{solution}
因为 \(3\mathbb Z+2^n\mathbb Z=\mathbb Z\)，闭包公式给出 \(\overline{3\mathbb Z}=\mathbb Z\)，故这个真子模不闭。局部化以后 \(3\) 成为单位，\(3\mathbb Z_{(2)}=\mathbb Z_{(2)}\)，当然闭。第一个环中 \((2)\not\subseteq J(\mathbb Z)\)，所以没有违反 Krull 交定理的闭性结论。
\end{solution}

\begin{exercise}\label{b3:ex:13-graded-not-ring}
设 \(k\) 为域，令
\[
R=k[x,y]/(x^2-y^3,xy,y^4),\qquad
S=k[x,y]/(x^2,xy,y^4).
\]
证明它们都是局部环，极大理想均由 \(x,y\) 的像生成，其伴随分次环同构；再证明 \(R,S\) 不同构。可比较平方为零的元素是否全部属于极大理想的平方。
\end{exercise}
\begin{solution}
两环都有 \(k\) 基 \(1,x,y,y^2,y^3\)：在 \(R\) 中用 \(x^2=y^3\) 消去 \(x^2\)，在 \(S\) 中用 \(x^2=0\)，其余混合项和 \(y^4\) 均为零。这些关系确实给出五维代数；结合性只须额外核对 \(x^2y=0=xyx\) 及 \(x^3=0\)。因此上述标准形唯一。两个理想 \(\mathfrak m=(x,y)\) 均满足 \(\mathfrak m^2=(y^2,y^3)\)、\(\mathfrak m^3=(y^3)\)、\(\mathfrak m^4=0\)，且商为 \(k\)，故是唯一极大理想。两环的逐层乘法都给出 \(k[X,Y]/(X^2,XY,Y^4)\)。

若 \(u\in R\) 满足 \(u^2=0\)，先模掉 \(\mathfrak m\) 得到常数项为零。写 \(u=ax+by+cy^2+dy^3\)，则
\[
u^2=b^2y^2+(a^2+2bc)y^3.
\]
由基的独立性先得 \(b=0\)，再得 \(a=0\)，所以 \(u\in\mathfrak m^2\)。但在 \(S\) 中，\(x^2=0\) 且 \(x\notin\mathfrak m^2\)。环同构保持极大理想及其平方，因而这两个环不同构。此例在任意特征下成立，说明即使底域和全部分次片均相同，分次环仍会遗失高阶乘法信息。
\end{solution}

\begin{exercise}\label{b3:ex:13-principal-rees}
设 \(R\) 为交换环，\(a\in R\) 不是零因子，证明 \(\mathcal R_{(a)}(R)\cong R[T]\)，其中 \(T\) 对应 \(at\)。再设 \(k\) 为域，若 \(R=k[\epsilon]/(\epsilon^2)\)、\(a=\epsilon\)，求此满射的核。
\end{exercise}
\begin{solution}
映射 \(R[T]\rightarrow R[t]\) 把 \(\sum r_iT^i\) 送到 \(\sum r_ia^it^i\)，像即 Rees 环。若 \(a\) 不是零因子，则每个 \(a^i\) 也不是零因子，逐系数比较得到核为零。第二种情形中次数至少二全部消失，次数一的系数须满足 \(r_1\epsilon=0\)，即 \(r_1\in(\epsilon)\)，常数项须为零。因此核为 \((\epsilon T,T^2)\)。
\end{solution}

\begin{exercise}\label{b3:ex:13-shifted-filtration}
设 \(R\) 为交换环，\(I\subseteq R\) 为理想，\(M\) 为 \(R\) 模，\(r\) 为非负整数。令 \(F^nM=I^{\max(n-r,0)}M\)。证明这是稳定 \(I\) 滤过；指出稳定开始的一个次数并比较其完备化。
\end{exercise}
\begin{solution}
下降性明显。当 \(n<r\) 时，\(IF^nM=IM\subseteq M=F^{n+1}M\)；当 \(n\ge r\) 时，\(IF^nM=I^{n-r+1}M=F^{n+1}M\)。所以可取稳定次数 \(r\)。滤过只把原滤过向后推迟有限步，与 \(I^nM\) 相互共尾，故给出相同拓扑和自然同构的完备化。
\end{solution}

\begin{exercise}\label{b3:ex:13-artin-rees-number}
设 \(k\) 为域，在 \(R=k[x,y]\) 中取 \(I=(x,y)\)、\(N=(x^a,y^b)\subseteq M=R\)，其中 \(a,b\ge1\)。求 Artin–Rees 等式中最小的 \(c\)。
\end{exercise}
\begin{solution}
令 \(c=\max(a,b)\)。单项式逐项判断给出
\[
N\cap I^n=x^aI^{n-a}+y^bI^{n-b}
=I^{n-c}(x^aI^{c-a}+y^bI^{c-b})\quad(n\ge c).
\]
所以此 \(c\) 可行。不妨 \(b\ge a\)。若 \(c'<b\)，取 \(n=b\)，则 \(y^b\in N\cap I^b\)，但它不属于 \(I^{b-c'}(N\cap I^{c'})\)：后一单项式理想的每个不被 \(x\) 整除的生成单项式，\(y\) 次数至少为 \(b+(b-c')>b\)。故最小值为 \(\max(a,b)\)。
\end{solution}

\begin{exercise}\label{b3:ex:13-krull-domain}
设 \(R\) 为 Noether 整环，\(I\) 为 \(R\) 的真理想。证明 \(\bigcap_nI^n=0\)，即使 \(I\) 不包含于 Jacobson 根也成立。
\end{exercise}
\begin{solution}
令 \(N=\bigcap_nI^n\)。它是 Noether 环 \(R\) 的理想，故有限生成；\cref{b3:thm:artin-rees}给出 \(N=IN\)。对 \(N\) 的恒等自同态应用\cref{b3:thm:determinant-trick}，得到某个 \(a\in I\) 使
\[
(1-a)N=0.
\]
因 \(I\ne R\)，有 \(1-a\ne0\)；整环中的非零元素不是零因子，所以 \(N=0\)。这里用整环的非零因子性质替代了 \(1-a\) 为单位的条件。
\end{solution}

\begin{exercise}\label{b3:ex:13-p-adic-digits}
设 \(p\) 为素数，\(\mathbb Z_p\) 为整数环的 \((p)\) 进完备化。证明每个 \(z\in\mathbb Z_p\) 唯一写成 \(\sum_{i\ge0}a_ip^i\)，其中各 \(a_i\) 为整数且 \(0\le a_i<p\)。证明 \(z\) 可逆当且仅当 \(a_0\ne0\)。
\end{exercise}
\begin{solution}
从第一个商选择唯一 \(a_0\)，再从第二个商中减去 \(a_0\)，剩余是唯一一个 \(a_1p\)；相容性使这种逐位选择持续且唯一。反向任意这样的数位族给出相容截断。如果 \(a_0=0\)，则 \(z\) 在剩余域中为零，不能可逆。若 \(a_0\ne0\)，先选整数 \(b\) 使 \(bz\equiv1\pmod p\)，再用 \((1-u)^{-1}=\sum_{j\ge0}u^j\)、\(u=1-bz\in p\mathbb Z_p\) 构造逆元。
\end{solution}

\begin{exercise}\label{b3:ex:13-composite-adic}
设 \(m=\prod_{i=1}^rp_i^{e_i}>1\)，其中 \(p_i\) 是不同素数。求 \(\mathbb Z\) 的 \((m)\) 进完备化。
\end{exercise}
\begin{solution}
中国剩余定理逐阶给出 \(\mathbb Z/m^n\mathbb Z\cong\prod_i\mathbb Z/p_i^{e_in}\mathbb Z\)，同构与过渡映射相容。逆极限与有限直积交换，因为相容条件逐坐标独立。指数列 \(e_in\) 在正整数中共尾，其极限与全部 \(p_i\) 次幂商的极限相同，所以完备化为 \(\prod_i\mathbb Z_{p_i}\)。
\end{solution}

\begin{exercise}\label{b3:ex:13-complete-finite-module}
设 \(R\) 为交换 Noether 环，\(I\subseteq R\) 为理想，且 \(R\) 对 \(I\) 进拓扑完备。证明每个有限生成的 \(R\) 模 \(M\) 都 \(I\) 进完备。取素数 \(p\)，用 \(R=\mathbb Z_p\)、\(I=(p)\)、\(M=\mathbb Q_p=\mathbb Z_p[1/p]\) 说明有限性不能省略。
\end{exercise}
\begin{solution}
\cref{b3:thm:completion-tensor}使自然映射等同于 \(M\rightarrow R\otimes_RM\)，因此是同构。另一方面，对 \(R=\mathbb Z_p\)、\(I=(p)\)、\(M=\mathbb Q_p=\mathbb Z_p[1/p]\)，有 \(p^nM=M\)，故 \(\widehat M=0\ne M\)。这里失败的首先是分离性。
\end{solution}

\begin{exercise}\label{b3:ex:13-completed-quotient}
设 \(k\) 为域。求 \(R=k[x,y]_{(x,y)}/(y^2,x^3)\) 的极大理想进完备化及其 \(k\) 维数。再求 \(k[x,y]_{(x,y)}/(y^2-x^3)\) 的极大理想进完备化。
\end{exercise}
\begin{solution}
第一环的极大理想四次幂为零，因为任何全次数至少四的单项式都被 \(x^3\) 或 \(y^2\) 整除，故它已完备。基为 \(1,x,x^2,y,xy,x^2y\)，维数六；与 \(k[\![x,y]\!]/(y^2,x^3)\) 相同。第二环由商完备化公式得到 \(k[\![x,y]\!]/(y^2-x^3)\)。这一结论不以把每个形式级数直接代入一个有理函数为依据，而来自全部有限阶商的相容同构。
\end{solution}

\begin{exercise}\label{b3:ex:13-completed-presentation}
设 \(k\) 为域，\(R=k[x]_{(x)}\)、\(A=\operatorname{diag}(x^2,x^4):R^2\rightarrow R^2\)、\(M=\operatorname{coker}A\)。求 \(M\) 的 \((x)\) 进完备化 \(\widehat M\) 的展示及长度。
\end{exercise}
\begin{solution}
在 \(\widehat R=k[\![x]\!]\) 上保留同一个矩阵，得
\(\widehat M\cong k[\![x]\!]/(x^2)\oplus k[\![x]\!]/(x^4)\)。两直和项分别有 \(k\) 基 \(1,x\) 与 \(1,x,x^2,x^3\)，故长度为六。原模已被 \(x^4\) 零化，各阶商最终恒定，所以自然映射 \(M\rightarrow\widehat M\) 本来就是同构。
\end{solution}

\begin{exercise}\label{b3:ex:13-idempotent-completion}
设 \(k\) 为域，取 \(R=k\times k\)、\(I=k\times0\)。求 \(I\) 进完备化映射，并直接证明它平坦而不忠实。
\end{exercise}
\begin{solution}
所有 \(R/I^n\) 都等于第二个 \(k\)，完备化就是投影 \((a,b)\mapsto b\)。作为 \(R\) 模，第二因子为幂等元 \((0,1)\) 生成的直和项，因此投射且平坦。非零模 \(k\times0\) 张量到第二因子后为零，所以不忠实。这也给出自然映射不单射的完备化例子。
\end{solution}

\begin{exercise}\label{b3:ex:13-finite-inverse}
设 \(k\) 为域。在 \(k[\![x,y]\!]/(x,y)^4\) 中求 \(1-x-y^2\) 的逆，写出所有全次数小于四的项。
\end{exercise}
\begin{solution}
令 \(h=x+y^2\)，所求为 \(1+h+h^2+h^3\) 的三阶截断。由于 \(h^2=x^2+2xy^2+y^4\)、\(h^3\equiv x^3\pmod{(x,y)^4}\)，答案是
\(1+x+y^2+x^2+2xy^2+x^3\)。系数 \(2\) 在特征二时为零；等比恒等式说明该式在任意特征都成立。
\end{solution}

\begin{exercise}\label{b3:ex:13-linear-equations-descend}
设 \((R,\mathfrak m)\) 为交换 Noether 局部环，\(\widehat R\) 为极大理想进完备化，\(A\) 是 \(R\) 上的 \(s\times r\) 矩阵，\(b\in R^s\)。证明：若 \(Ax=b\) 在 \(\widehat R^r\) 中有解，则在 \(R^r\) 中也有解。说明这个结论为何不保证选定的形式解本身属于 \(R^r\)。
\end{exercise}
\begin{solution}
令 \(N=\operatorname{im}(A)\subseteq R^s\)。完备化正合性说明 \(\operatorname{im}(\widehat A)=\widehat N\) 在 \(\widehat R^s\) 中的像。形式解存在意味着 \(b\) 的像属于它，由\cref{b3:prop:completion-ideal-contraction}得 \(b\in N\)，即原环中有解。解的选择可能相差核中的任意元素；例如零方程 \(0x=0\) 的任意形式级数都是解，却不必来自 \(R\)。
\end{solution}

\begin{exercise}\label{b3:ex:13-cusp-multiplicity}
设 \(k\) 为域，\(R=k[x,y]_{(x,y)}/(y^2-x^3)\)，\(B=\widehat R\) 为极大理想进完备化。求两环的维数和极大理想进重数。
\end{exercise}
\begin{solution}
\cref{b3:prop:hypersurface-tangent-cone}给出伴随分次环 \(k[X,Y]/(Y^2)\)，其次数零维数一，每个正次数维数二。故累计长度为 \(2n+1\)。由 Hilbert–Samuel 定理，\(\dim R=1\)、\(e(R)=2\)。完备化各阶商相同，得到 \(\dim B=1\)、\(e(B)=2\)。重数是此指定极大理想的重数。
\end{solution}

\begin{exercise}\label{b3:ex:13-approximate-membership}
设 \((R,\mathfrak m)\) 为交换 Noether 局部环，\(J\subseteq R\) 为理想，\(f\in R\)。证明 \(f\in J\) 当且仅当对每个整数 \(N\ge1\) 有 \(f\in J+\mathfrak m^N\)。给出只检查某个固定 \(N\) 不足的例子。
\end{exercise}
\begin{solution}
反向在有限生成模 \(R/J\) 中使用\cref{b3:thm:krull-intersection}，正向显然。设 \(k\) 为域。对任意固定 \(N\)，在 \(R=k[x]_{(x)}\) 中取 \(J=(x^{N+1})\)、\(f=x^N\)。则 \(f\in J+(x)^N\)，却 \(f\notin J\)。所以量词“每个 \(N\)”不能替换成一个事先任意选定的精度。
\end{solution}

\begin{exercise}\label{b3:ex:13-formal-nonlinear-root}
设 \(k\) 为域且 \(\operatorname{char}k\ne2\)。在 \(k[\![x]\!]\) 中构造常数项为一、满足 \(u^2=1+x\) 的唯一级数，并说明在 \(k=\mathbb Q\) 时它不属于 \(\mathbb Q[x]_{(x)}\)。
\end{exercise}
\begin{solution}
写 \(u=1+\sum_{n\ge1}a_nx^n\)。比较系数得 \(2a_1=1\)，且对 \(n\ge2\)，
\(2a_n=-\sum_{i=1}^{n-1}a_ia_{n-i}\)。由于 \(2\) 可逆，递归唯一决定所有系数，并保证每阶方程，故给出唯一形式解。若它在 \(\mathbb Q(x)\) 中，写成互素多项式比 \(P/Q\)，则 \(P^2=(1+x)Q^2\)。不可约因子 \(1+x\) 在左侧出现偶数次，右侧出现奇数次，矛盾。于是形式解的存在并不一般下降为原局部环中非线性方程的解。
\end{solution}

\begin{optionalexercise}\label{b3:ex:13-nonextended-ideal}
设 \(k\) 为可数域。证明存在 \(f\in xk[\![x]\!]\) 在 \(k(x)\) 上超越；再证明 \(k[\![x,y]\!]\) 中的非零理想 \(\bigl(y-f(x)\bigr)\) 与 \(k[x,y]_{(x,y)}\) 的交为零。
\end{optionalexercise}
\begin{solution}
\(k[x,T]\) 可数，每个非零多项式在域 \(k((x))\) 中只有有限个根，因此在 \(k(x)\) 上代数的级数只有可数多个。所有系数仅取两个不同值的级数已不可数，故可选超越 \(f\)，减去其常数项后仍超越。

写 \(F(x,y)=\sum_{i,j\ge0}a_{ij}x^iy^j\)。由于 \(f(x)\in xk[\![x]\!]\)，项 \(x^if(x)^j\) 的阶数至少为 \(i+j\)。在任一固定精度下，代入和因而只涉及有限多个 \((i,j)\)，所以 \(F(x,f(x))\) 良定义，并给出连续满射 \(k[\![x,y]\!]\rightarrow k[\![x]\!]\)。利用有限恒等式
\[
y^j-f(x)^j=(y-f(x))\sum_{\ell=0}^{j-1}y^{j-1-\ell}f(x)^\ell
\quad(j\ge1),
\]
构造
\[
G(x,y)=\sum_{i\ge0,\,j\ge1}a_{ij}x^i
       \sum_{\ell=0}^{j-1}y^{j-1-\ell}f(x)^\ell.
\]
其中每个加项的全次数至少为 \(i+j-1\)，故任一给定总次数只收到有限多个 \((i,j,\ell)\) 的贡献。于是 \(G\) 也是良定义的形式级数，逐有限阶使用上述恒等式得到
\[
F(x,y)-F(x,f(x))=(y-f(x))G(x,y).
\]
这证明代入映射的核恰为 \(\bigl(y-f(x)\bigr)\)。若原局部环的 \(P/Q\) 落入此核，\(Q(0,0)\ne0\) 使 \(Q(x,f(x))\) 可逆，因此 \(P(x,f(x))=0\)。超越性迫使 \(P=0\)。故该非零理想不是收缩理想的扩张，说明并非完备环中的所有理想都来自原环。
\end{solution}
